\documentclass[aspectratio=169]{beamer}

\usepackage{tikz}
\usepackage{styles/slidesmainstyle10}

% ============================================================
% METADADOS DA AULA
% ============================================================
\newcommand{\LessonTitle}{Geração de Dados Sintéticos de Qualidade}
\newcommand{\LessonBg}{img/bg.png}
\newcommand{\LessonAuthor}{Prof. Jeff Oliveira, Ph.D.}
\newcommand{\LessonInfo}{Caso IPT | SI10 | 2026}

\begin{document}

\coverslide{\LessonTitle}{\LessonBg}{\LessonAuthor}{\LessonInfo}

\textslide
  {MAPA DA AULA}
  {}
  {
    Vamos acompanhar o mesmo caminho do notebook, mas em versão enxuta:

    \vspace{0.7em}
    \textbf{Passo 0:} fixar semente e garantir reprodutibilidade.\\
    \textbf{Passos 1, 2A e 2B:} entender o caso, identificar lacunas e propor hipóteses.\\
    \textbf{Passos 3 e 4:} decidir a unidade de análise e o schema da base.\\
    \textbf{Passos 5, 6, 7 e 8:} escolher distribuições, gerar sessões e inserir outliers plausíveis.

    \vspace{0.8em}
    A meta não é sofisticação. A meta é produzir uma base sintética simples e justificável.
  }

\sectionslide{ENTENDIMENTO DO PROBLEMA}

\textslide
  {O CASO IPT}
  {}
  {
    O IPT quer tomar decisões sobre comportamento de navegação na intranet.
    O problema é que os dados disponíveis são agregados e não mostram diretamente
    o comportamento de cada visita.

    \vspace{0.8em}
    A aula mostra como sair desse ponto e montar uma base sintética mínima,
    coerente com o caso e fácil de explicar.
  }

\bulletslide
  {O QUE O PARCEIRO QUER}
  {}
  {
    \item Reduzir o tempo gasto procurando links operacionais
    \item Aumentar o engajamento com comunicados internos
    \item Segmentar usuários por perfil de comportamento
    \item Sustentar decisões de redesign e testes A/B
  }

\questionslide{%
  \ql{O que sabemos do caso?}\\[0.3em]
  \qh{Quais sinais já aparecem}\\[0.3em]
  \ql{mesmo com dados agregados?}
}

\stmtslide
  {Para manter a aula simples}
  {vamos trabalhar um único objetivo de negócio}

\textslide
  {RECORTE DA AULA}
  {}
  {
    Nesta aula, vamos focar em:

    \vspace{0.6em}
    \textbf{aumentar o engajamento com comunicados internos}

    \vspace{0.8em}
    Esse recorte reduz a complexidade, mas preserva o raciocínio central
    de geração de dados sintéticos.
  }

\textslide
  {O QUE A PLANILHA CONTÉM}
  {}
  {
    \renewcommand{\arraystretch}{1.2}
    \small
    \begin{tabular}{p{3.1cm} p{4.9cm} p{2.5cm}}
    \textbf{Aba} & \textbf{O que contém} & \textbf{Granularidade} \\
    \\
    Tráfego geral & Visualizadores e visitas por dia & 1 linha = 1 dia \\
    Conteúdo popular & Itens mais visitados nos últimos 7 dias & 1 linha = 1 item \\
    Uso por dispositivo & Desktop, mobile app, mobile web, tablet & 1 linha = 1 dia \\
    Uso por tempo & Média de viewers por dia e hora & 1 linha = 1 janela \\
    \end{tabular}
  }

\bulletslide
  {O QUE JÁ CONSEGUIMOS AFIRMAR}
  {}
  {
    \item O acesso é majoritariamente desktop
    \item Há janelas claras de pico entre 8h e 11h
    \item Alguns conteúdos concentram grande parte das visitas
    \item O dado é agregado, não comportamental por sessão
  }

\questionslide{%
  \ql{O que falta nos dados?}\\[0.3em]
  \qh{Onde estão as lacunas}\\[0.3em]
  \ql{que impedem análise comportamental direta?}
}

\bulletslide
  {LACUNAS PRINCIPAIS}
  {}
  {
    \item Não há sessão individual
    \item Não há perfil de usuário observável
    \item Não há sequência de navegação até o comunicado
    \item Não há sinal claro do que aconteceu depois da leitura
  }

\questionslide{%
  \ql{Como virar lacuna em hipótese?}\\[0.3em]
  \qh{Se o dado não existe,}\\[0.3em]
  \ql{o que podemos assumir com critério?}
}

\textslide
  {DA LACUNA À HIPÓTESE}
  {}
  {
    \textbf{Lacuna:} não sabemos quem lê comunicados, nem como a leitura
    se distribui ao longo das sessões.

    \vspace{0.8em}
    \textbf{Hipótese simples:} em dias úteis, a maior parte das sessões de
    leitura ocorre em desktop, entre 8h e 11h, com alguns perfis lendo mais
    do que outros.
  }

\questionslide{%
  \ql{O que cada linha da base representa?}\\[0.3em]
  \qh{Usuário, sessão, página}\\[0.3em]
  \ql{ou clique?}
}

\stmtslide
  {Decisão de modelagem}
  {1 linha = 1 sessão}

\sectionslide{MODELAGEM DA BASE}

\questionslide{%
  \ql{Quais colunas a base terá?}\\[0.3em]
  \qh{Antes de gerar,}\\[0.3em]
  \ql{precisamos nomear as variáveis.}
}

\begin{singlecodeslide}{SCHEMA MINIMO DA BASE}{Python}
\begin{lstlisting}
colunas = [
    "session_id",
    "perfil_usuario",
    "tipo_dispositivo",
    "dia_semana",
    "faixa_horaria",
    "origem_acesso",
    "objetivo_sessao",
    "pageviews",
    "tempo_busca_segundos",
    "tempo_leitura_segundos",
    "tempo_total_segundos",
    "usou_busca",
    "clicou_comunicado",
    "sucesso_encontrou_o_que_queria",
    "abandono_rapido",
]
\end{lstlisting}
\end{singlecodeslide}

\questionslide{%
  \ql{Qual distribuicao usar?}\\[0.3em]
  \qh{A forma estatistica depende}\\[0.3em]
  \ql{do tipo de comportamento modelado.}
}

\rightimgslide
  {BERNOULLI}
  {Use Bernoulli quando a variavel so pode assumir dois valores.\\[0.7em]
   Exemplo trivial: \texttt{clicou\_comunicado = 1} ou \texttt{0}.}
  {img/dist_bernoulli.png}
  {0.48}
  {true}

\rightimgslide
  {LOG-NORMAL}
  {Use log-normal quando a variavel e sempre positiva e assimetrica.\\[0.7em]
   Exemplo trivial: \texttt{tempo\_leitura\_segundos}.}
  {img/dist_lognormal.png}
  {0.48}
  {true}

\rightimgslide
  {POISSON}
  {Use Poisson quando a variavel representa contagens em um intervalo.\\[0.7em]
   Exemplo trivial: \texttt{pageviews} na sessao.}
  {img/dist_poisson.png}
  {0.48}
  {true}

\begin{singlecodeslide}{REGRAS DE GERACAO}{Python}
\begin{lstlisting}
def escolher_objetivo(perfil):
    if perfil == "pesquisador":
        return sorteio({
            "ler_comunicado": 0.45,
            "buscar_documento": 0.35,
            "buscar_sistema": 0.20,
        })
    return sorteio({
        "ler_comunicado": 0.25,
        "buscar_documento": 0.25,
        "buscar_sistema": 0.50,
    })

def prob_clicar_comunicado(objetivo, faixa_horaria):
    p = 0.12
    if objetivo == "ler_comunicado":
        p += 0.58
    if faixa_horaria == "08-11h":
        p += 0.08
    return min(p, 0.95)
\end{lstlisting}
\end{singlecodeslide}

\begin{singlecodeslide}{GERANDO UMA SESSAO}{Python}
\begin{lstlisting}
for session_id in range(n):
    perfil = escolher_perfil()
    dispositivo = escolher_dispositivo()
    dia = escolher_dia_semana()
    faixa = escolher_faixa_horaria()
    objetivo = escolher_objetivo(perfil)

    pageviews = np.random.poisson(lam_por_objetivo[objetivo])
    tempo_busca = np.random.lognormal(*param_busca[objetivo])
    tempo_leitura = np.random.lognormal(*param_leitura[objetivo])
    clicou = bernoulli(prob_clicar_comunicado(objetivo, faixa))

    linha = montar_sessao(
        session_id, perfil, dispositivo, dia, faixa,
        objetivo, pageviews, tempo_busca, tempo_leitura, clicou
    )
    linhas.append(linha)
\end{lstlisting}
\end{singlecodeslide}

\sectionslide{REALISMO E FECHAMENTO}

\begin{singlecodeslide}{INSERINDO OUTLIERS}{Python}
\begin{lstlisting}
idx_tempo = df.sample(frac=0.006, random_state=42).index
df.loc[idx_tempo, "tempo_leitura_segundos"] *= np.random.uniform(2.2, 4.0, len(idx_tempo))

idx_pv = df.sample(frac=0.004, random_state=7).index
df.loc[idx_pv, "pageviews"] += np.random.randint(8, 18, len(idx_pv))

df["abandono_rapido"] = (
    (df["tempo_total_segundos"] < 35) &
    (df["pageviews"] <= 2)
).astype(int)
\end{lstlisting}
\end{singlecodeslide}

\bulletslide
  {O QUE LEVAR DESTA AULA}
  {}
  {
    \item Fixar semente ajuda a reproduzir o experimento
    \item Lacuna e hipotese nao sao a mesma coisa
    \item Unidade de analise define o que a base consegue mostrar
    \item Dado sintetico bom e dado simples, explicavel e justificavel
  }

\finalslide{img/final-img.png}

\end{document}
